\chapter{\texttt{medium.solar} --- adiabatic MSW propagation in the Sun}
\label{ch:solar}

Solar neutrinos are produced deep in the high-density solar core and exit
through a smoothly, monotonically decreasing electron density to vacuum.
This chapter's physics is dominated by one fact exploited throughout: the
density gradient is slow enough, on the scale of the relevant MSW
resonances, that the \textbf{adiabatic approximation} applies almost
everywhere in the standard oscillation parameter space -- a neutrino
produced in a given local matter eigenstate stays in it as the density
drops, so propagation reduces to an incoherent mixture of eigenstates
rather than a coherent evolutor (see the note in
\cref{sec:solar-no-evolutor}).

\section{\texttt{profile.py} --- the tabulated solar model}
\label{sec:solar-profile}

\modulehead{medium/solar/profile.py}{\pyclass{SolarProfile}: tabulated
solar radius, electron density, per-source production fractions, and
per-source total fluxes.}

\begin{implbox}
\pyclass{SolarProfile} stores four tabulated quantities on a common radius
grid $r/R_\odot\in[0,1]$: the electron density $n_e(r)$ (mol/cm$^3$), a
production-fraction profile $f_{\rm source}(r)$ per neutrino-producing
reaction (pp, $^8$B, $^7$Be, ...), and a total integrated flux
$\Phi_{\rm source}$ per source. \pyfunc{SolarProfile.default} loads these
from the package's bundled CSV tables (\code{tpeanuts/data/solar}), whose
paths are fixed by \code{tpeanuts.config.default.solar\_model\_filename}
and \code{solar\_fluxes\_filename} -- currently the \textbf{zenodo
SF3-AGSS09} tables, a standard-solar-model dataset combining helioseismic
constraints with the AGSS09 heavy-element abundances
\parencite{Vinyoles2017B16}, tabulated over the full solar radius (2001
rows). \pyclass{SolarParameters} allows overriding both file paths
explicitly, but the defaults are the project's validated reference
composition and should not be changed casually.
\begin{itemize}
  \item \pyfunc{electron\_density(r\_query)} / \pyfunc{production\_fraction(source, r\_query)}:
    linear interpolation on the tabulated grid (\pyfunc{util.math.interp1d\_linear}),
    clamped to the boundary values outside $[0,1]$.
  \item \pyfunc{source\_fractions(sources)}: stacks one or several tabulated
    (not interpolated) production-fraction profiles.
  \item \pyfunc{normalized\_fraction(source)}: normalises a production
    profile to unit trapezoidal integral over $r$
    (\pyfunc{util.math.normalize\_trapz}).
  \item \pyfunc{flux(source)}: the total (radius-integrated) flux
    normalisation for one source, as tabulated.
  \item \code{use\_LZ}: boolean flag (default \code{False}) selecting
    whether \texttt{probability.py} (\cref{sec:solar-probability}) applies
    the Landau-Zener correction of \cref{sec:solar-lz}.
\end{itemize}
\end{implbox}

\begin{notebox}
\textbf{Validity of the tabulated model.} The bundled solar model is a
\emph{standard} solar model: a hydrostatic, helioseismically calibrated
stellar-structure calculation, not a direct measurement. Its accuracy is
therefore tied to the input microphysics (opacities, nuclear reaction
rates, heavy-element abundances) and is typically quoted at the few-percent
level for the core density and composition profiles that set $n_e(r)$ and
the production-fraction shapes; solar-neutrino and helioseismic data
provide independent, largely consistent cross-checks of this precision.
Because $n_e(r)$ only ever enters through the dimensionless ratio $V_k$
(\cref{sec:solar-matter-mixing}), a uniform rescaling error in the model's
overall density normalisation would bias the resonance radius, not the
qualitative adiabatic/non-adiabatic character of the propagation.
\end{notebox}

\section{\texttt{matter\_mixing.py} --- the adiabatic matter-mixing angles}
\label{sec:solar-matter-mixing}

\modulehead{medium/solar/matter\_mixing.py}{Closed-form matter-modified
mixing angles $\theta_{12}^M(r)$, $\theta_{13}^M(r)$ used by the adiabatic
solar-probability approximation.}

\begin{theorybox}
Because $\Delta m^2_{3\ell}\gg\Delta m^2_{21}$ and $\theta_{13}$ is small,
the solar MSW problem separates hierarchically into two sequential
effective 2-flavour resonances (the $1$--$3$ sector, then the $1$--$2$
sector), each with the standard closed-form matter-mixing-angle solution
\parencite{GiuntiKim2007} -- no numerical diagonalisation of the full
$3\times3$ Hamiltonian is needed in the Standard Model case. Define the
dimensionless matter-potential ratio and the effective atmospheric
splitting
\begin{align}
  V_k(\Delta m^2,E,n_e) &= \frac{V(n_e)}{k(\Delta m^2,E)}
  = \frac{2E\hbar c\,V_{CC}}{\Delta m^2}, \\
  \Delta m^2_{ee} &= \cos^2\theta_{12}\,\Delta m^2_{31} + \sin^2\theta_{12}\,\Delta m^2_{32},
\end{align}
using $V$ and $k$ exactly as defined in \cref{ch:core-common}
(\texttt{potential.py} section) -- $V_k$ is simply their ratio, evaluated
here for a given splitting rather than embedded in the Hamiltonian. The
matter-modified reactor angle is
\begin{equation}
  \keyresult{\theta_{13}^M = \tfrac12\arccos\!\left(
    \frac{\cos2\theta_{13}-V_k(\Delta m^2_{ee},E,n_e)}
         {\sqrt{\bigl(\cos2\theta_{13}-V_k(\Delta m^2_{ee},E,n_e)\bigr)^2+\sin^22\theta_{13}}}
  \right)},
\end{equation}
and, using $\theta_{13}^M$ as an input, the matter-modified solar angle is
\begin{equation}
  \keyresult{\theta_{12}^M = \tfrac12\arccos\!\left(
    \frac{\cos2\theta_{12}-V_k'}
         {\sqrt{(\cos2\theta_{12}-V_k')^2+\sin^22\theta_{12}\cos^2(\theta_{13}^M-\theta_{13})}}
  \right)},
\end{equation}
\begin{equation}
  V_k' = V_k(\Delta m^2_{21},E,n_e)\cos^2\theta_{13}^M
       + \frac{\Delta m^2_{ee}}{\Delta m^2_{21}}\sin^2(\theta_{13}^M-\theta_{13}).
\end{equation}
Each resonance occurs where the corresponding $V_k=\cos2\theta$, i.e.\
where the vacuum mixing angle is driven to $\pi/4$ by the matter
potential; away from resonance $\theta^M\to\theta$ (low density) or
$\theta^M\to\pi/2$ (high density), recovering the vacuum and
fully-suppressed limits respectively.
\end{theorybox}

\begin{notebox}
\code{legacy\_precision=True} reproduces the legacy \texttt{peanuts}
combined prefactor for $V_k$ ($3.868\times10^{-7}/2.533$, both factors
independently rounded to 4 significant digits in the original
implementation) rather than deriving the matter and kinetic prefactors
separately from full-precision constants (\cref{ch:core-common}); the two
choices differ by a relative $\sim3.6\times10^{-4}$ -- an order of
magnitude larger than the matter-potential-only rounding error of
\cref{ch:core-common}, since both legacy prefactors compound. Intended
purely for bit-comparable validation against legacy \texttt{peanuts}.
\end{notebox}

\begin{implbox}
\pyfunc{Vk(Deltam2, E, ne, antinu, legacy\_precision)},
\pyfunc{DeltamSqee(oscillation)}, \pyfunc{th13\_M(oscillation, E, ne, ...)},
\pyfunc{th12\_M(oscillation, E, ne, ..., th13m=None)} implement the four
formulas above directly; \pyfunc{th12\_M} accepts a precomputed
\code{th13m} to avoid recomputing it when both angles are needed together
(as \pyfunc{Tei} below does).
\end{implbox}

\section{\texttt{landau\_zener.py} --- non-adiabatic corrections}
\label{sec:solar-lz}

\modulehead{medium/solar/landau\_zener.py}{Landau-Zener transition
probability $P_{LZ}(E)$ quantifying the breakdown of the adiabatic
approximation near the MSW resonance.}

\begin{theorybox}
At \emph{finite} adiabaticity, a neutrino crossing the MSW resonance can
jump between the two local matter eigenstates instead of adiabatically
following one of them. The exponential Landau-Zener approximation
\parencite{Parke1986} gives the jump probability as
\begin{equation}
  \keyresult{P_{LZ} = \exp\!\left(-\frac{\pi}{2}\gamma_{\rm res}\right)},
\end{equation}
with the adiabaticity parameter evaluated at the resonance
\parencite{GiuntiKim2007}
\begin{equation}
  \gamma_{\rm res} = \frac{\Delta m^2_{21}\sin^22\theta_{12}}{2E\hbar c\cos2\theta_{12}}\,L_n(r_{\rm res}),
  \qquad
  L_n = \left|\frac{n_e}{\mathrm{d}n_e/\mathrm{d}l}\right|_{\rm res},
\end{equation}
where $L_n$ is the electron-density scale height (in metres) at the
resonance radius $r_{\rm res}(E)$, defined by $V_k(\Delta m^2_{21},E,n_e(r_{\rm
res}))=\cos2\theta_{12}$ (i.e.\ where $\theta_{12}^M=\pi/4$). Large
$\gamma_{\rm res}$ means a slowly varying density on the scale of the
oscillation length at resonance, hence a vanishing jump probability
(adiabatic); small $\gamma_{\rm res}$ means a comparatively abrupt density
change, hence a sizeable jump probability.
\end{theorybox}

\begin{notebox}
\textbf{Scope and limitations, as documented in the module itself.} Only
the $\theta_{12}$-sector resonance is tracked ($\theta_{13}$'s resonance
sits deep in the solar core at very high density and is negligible for
every standard solar source). $P_{LZ}=0$ is returned wherever no resonance
exists in the solar volume: below-threshold energies, antineutrinos (no MSW
resonance for standard parameters), and LMA-Dark parameters
($\theta_{12}>\pi/4$, resonance at unphysical negative density). For
standard LMA parameters and pp/$^8$B solar neutrinos,
$\gamma_{\rm res}\sim10^3$--$10^4$, so $P_{LZ}\sim e^{-1500}\approx0$ and
the pure adiabatic approximation is excellent; this module becomes
relevant mainly for non-standard parameter explorations or precision
studies in the 1--3 MeV range. Solar NSI is \textbf{not} included here: the
resonance condition uses the standard $V_k$, and \pyfunc{Tei}
(\cref{sec:solar-probability}) silently ignores \code{p\_lz} once
\code{oscillation.nsi} is set, since both the resonance condition and the
adiabaticity parameter would need to be rederived for off-diagonal
$\epsilon$.
\end{notebox}

\begin{implbox}
\pyfunc{density\_gradient(solar\_profile)}: $\mathrm{d}n_e/\mathrm{d}\hat r$
on the tabulated grid via central differences (one-sided at the
boundaries). \pyfunc{resonance\_radius(oscillation, E, solar\_profile)}:
locates $r_{\rm res}(E)$ by linear interpolation at the first
sign-change of $V_k-\cos2\theta_{12}$ along the (monotonically decreasing)
density profile; returns \code{NaN} where no crossing exists.
\pyfunc{plz(oscillation, E, solar\_profile)}: the $P_{LZ}(E)$ formula
above, masked to exactly zero wherever \pyfunc{resonance\_radius} returns
\code{NaN}.
\end{implbox}

\begin{notebox}[title={No Coherent Medium Evolutor}]
\label{sec:solar-no-evolutor}
Unlike every other medium in this document, \texttt{medium.solar} has
\textbf{no \texttt{evolutor.py}}. This is a direct consequence of the
adiabatic approximation, not an omission: once a neutrino is assigned to a
local matter eigenstate at its production radius (\pyfunc{Tei}, next
section), the adiabatic theorem guarantees it stays in the
\emph{corresponding} eigenstate all the way to the solar surface, where the
density has dropped to (effectively) zero and the local matter eigenstates
have smoothly become the \emph{vacuum} mass eigenstates. No further
coherent phase evolution needs computing: the production-point projection
already \emph{is} the final incoherent mass-basis composition, and only the
final flavour projection (\pyfunc{solar\_probability\_state}) remains.
Consequently \texttt{medium.solar} exposes no
\pyfunc{solar\_probability\_transition}: unlike the other media
(\cref{ch:vacuum,ch:earth,ch:atmosphere}), there is no coherent evolutor $S$
to build a transition matrix from, so \pyfunc{solar\_probability\_state} is
the lowest-tier public probability function for this medium.
\end{notebox}

\section{\texttt{probability.py} --- production weights and flavour probabilities}
\label{sec:solar-probability}

\modulehead{medium/solar/probability.py}{\pyfunc{Tei} (production-point
matter-eigenstate weights), \pyfunc{solar\_\allowbreak{}probability\_\allowbreak{}mass} (integrated
over the production profile), \pyfunc{solar\_\allowbreak{}probability\_\allowbreak{}state} (final
flavour probabilities), and \pyfunc{solar\_\allowbreak{}probability\_\allowbreak{}integrated}
(spectrum-weighted energy average).}

\begin{theorybox}
\subsection*{\pyfunc{Tei}: flavour-to-matter-basis production weights}
$T_{ei}$ answers: given an electron neutrino produced at electron density
$n_e$ and energy $E$, what is the probability it is projected onto local
matter eigenstate $i$? In the Standard Model (adiabatic, no LZ) case this
is read directly off the matter-mixing angles of
\cref{sec:solar-matter-mixing}:
\begin{equation}
  \keyresult{
  \bigl(w_1,w_2,w_3\bigr) =
  \Bigl(\cos^2\theta_{13}^M\cos^2\theta_{12}^M,\ \
        \cos^2\theta_{13}^M\sin^2\theta_{12}^M,\ \
        \sin^2\theta_{13}^M\Bigr)
  },
\end{equation}
i.e.\ $|\langle\nu_e|\nu_i^M\rangle|^2$ for the standard PMNS-like
parametrisation applied to the matter-modified angles. When a
Landau-Zener probability $P_{LZ}$ is supplied (\cref{sec:solar-lz}), the
$1$--$2$ sector weights mix according to the LZ jump probability,
\begin{equation}
  w_1 = (1-P_{LZ})\cos^2\theta_{12}^M + P_{LZ}\sin^2\theta_{12}^M,
  \qquad
  w_2 = (1-P_{LZ})\sin^2\theta_{12}^M + P_{LZ}\cos^2\theta_{12}^M,
\end{equation}
while $w_3=\sin^2\theta_{13}^M$ is untouched (the $\theta_{13}$ resonance
is negligible, \cref{sec:solar-lz}). \textbf{When NSI is active}
(\code{oscillation.nsi} set), this entire analytic construction is
replaced by direct numerical diagonalisation: the full flavour-basis
Hamiltonian $H=H_{\rm kin}+H_{\rm mat}^{\rm NSI}$
(\pyfunc{hamiltonian\_flavour}, \cref{ch:core-common,ch:core-bsm}) is built
at each $(E,n_e)$ point and diagonalised with \pyfunc{torch.linalg.eigh};
the production weights are then simply the squared electron-flavour
components of each matter eigenvector,
\begin{equation}
  w_i(E,n_e) = \bigl|\langle\nu_e|\nu_i^M(E,n_e)\rangle\bigr|^2 = |{\rm eigvec}[\dots,0,i]|^2,
\end{equation}
which supports arbitrary off-diagonal $\epsilon$ (including LMA-Dark-type
degenerate scenarios) at the cost of losing the closed form; \code{p\_lz}
is silently ignored in this path.

\subsection*{From production weights to a final flavour probability}
\pyfunc{solar\_probability\_mass} integrates $T_{ei}(E,n_e(r))$ over the
production-radius distribution of one or more sources, weighted by each
source's normalised production-fraction profile $f_{\rm source}(r)$
(\cref{sec:solar-profile}):
\begin{equation}
  \keyresult{P_i(E) = \frac{\int_0^1 T_{ei}\bigl(E,n_e(r)\bigr)\,f_{\rm source}(r)\,\mathrm{d}r}
                          {\int_0^1 f_{\rm source}(r)\,\mathrm{d}r}},
\end{equation}
evaluated by the trapezoidal rule, giving an energy-dependent (not
radius-resolved) incoherent \emph{vacuum} mass-basis probability vector
$P_i(E)$ for that source -- by the argument of the note above, this
average over production radius already \emph{is} the mass-basis
composition reaching the solar surface, with no further coherent step
required. \pyfunc{solar\_probability\_state} performs the last, purely
kinematic conversion to flavour, using the ordinary incoherent-mixture
formula of
\cref{ch:core-common} (\texttt{probability.py} section) with the
\emph{identity} evolutor (there is nothing left to evolve coherently):
\begin{equation}
  \keyresult{P_\alpha(E) = \sum_i |U_{\alpha i}|^2\,P_i(E)},
\end{equation}
where $U$ is the ordinary \emph{vacuum} PMNS matrix
(\cref{ch:core-common,ch:core-sm}) -- the matter-modified angles only ever
appear inside $T_{ei}$, at the production point; every subsequent step in
this chapter works with the vacuum mixing matrix.
\end{theorybox}

\begin{implbox}
\begin{itemize}
  \item \pyfunc{Tei(oscillation, E, ne, p\_lz=None, legacy\_precision=False)}:
    the weight formulas above.
  \item \pyfunc{solar\_probability\_mass(oscillation, E, profile, sources, ...)}:
    builds the spatially resolved LZ mask (production radii above the
    resonance density only, \cref{sec:solar-lz}) when
    \code{profile.use\_LZ} and no NSI is active, then integrates $T_{ei}$
    against the source's production-fraction profile as above.
  \item \pyfunc{solar\_probability\_state(oscillation, E, profile, sources, ...)}:
    the final flavour-probability formula above, via
    \pyfunc{core.common.probability.probability\_incoherent}
    (\cref{ch:core-common}) with the identity operator standing in for
    ``no further coherent evolution''.
  \item \pyfunc{solar\_probability\_integrated(oscillation, E, profile, sources, spectrum, ...)}:
    builds \pyfunc{solar\_probability\_state} on an energy grid and
    averages it with
    \pyfunc{core.common.probability.probability\_integrated}
    (\cref{ch:core-common}), weighted by an explicit production
    \code{spectrum}.
\end{itemize}
\end{implbox}

\section{\texttt{flux.py} --- solar flux}

\modulehead{medium/solar/flux.py}{\pyfunc{solar\_\allowbreak{}flux\_\allowbreak{}state} and
\pyfunc{solar\_\allowbreak{}flux\_\allowbreak{}integrated}: apply per-source flux normalisations,
optional spectra, and energy integration to
\pyfunc{solar\_\allowbreak{}probability\_\allowbreak{}state}'s output.}

\begin{implbox}
\pyfunc{solar\_flux\_state(sources, profile, oscillation, E\_MeV, source\_spectrum=None, ...)}
adds no new physics: it calls \pyfunc{solar\_probability\_state} above for
the requested source(s), then \pyfunc{core.common.flux.flux\_state}
(\cref{ch:core-common}) to apply each source's total flux normalisation
(\pyfunc{SolarProfile.flux}) and optional energy spectrum -- the same
minimal pattern established in \cref{ch:vacuum}.
\pyfunc{solar\_flux\_integrated(...)} builds \pyfunc{solar\_flux\_state} on
an energy grid and integrates it with
\pyfunc{core.common.flux.flux\_integrated}, returning an unnormalised
physical rate.
\end{implbox}
